\chapter{Images, Convolution, and the DFT}
\label{ch:images-convolution-dft}

\begin{draftbox}
This chapter is scaffolded. It will absorb the DFT material and image
processing lab ideas.
\end{draftbox}

\section{Guiding Question}

Why does the Fourier basis make convolution simple?

\section{Planned Core Ideas}

\begin{itemize}
  \item images and signals can be treated as vectors;
  \item convolution is a linear map;
  \item circulant matrices are diagonalized by Fourier vectors;
  \item the DFT converts convolution into multiplication;
  \item linear filters can blur, sharpen, or detect changes.
\end{itemize}

\section{Planned Python Thread}

Implement tiny circular convolution, compare with FFT-based multiplication, and
apply small filters to arrays before moving to full image labs.

\section*{Chapter Summary}
\addcontentsline{toc}{section}{Chapter Summary}

To be completed in the chapter-shaping pass.

\section*{Exercises}
\addcontentsline{toc}{section}{Exercises}
\subsection*{A. Theory}
\begin{exercise}\exerciseid{18.A1} Explain why convolution with a fixed kernel is a linear map. \end{exercise}
\subsection*{B. By Hand}
\begin{exercise}\exerciseid{18.B1} Build the $4\times 4$ circulant matrix corresponding to the kernel $(1,0,-1,0)$. \end{exercise}
\subsection*{C. Python / NumPy}
\begin{exercise}\exerciseid{18.C1} Implement circular convolution directly for vectors of length $N$ and compare with multiplication by a circulant matrix. \end{exercise}
\subsection*{D. Challenge}
\begin{exercise}\exerciseid{18.D1} Explain why diagonalizing convolution is useful for fast computation. \end{exercise}

